
\begin{proposition}
\label{V.5.8}
El funtor $X\mto E_X$ es una equivalencia de la categoría~$\cal{C}$ con la
categoría de los funtores covariantes del grupoide fundamental~$\Gamma$
de~$\cal{C}$ en~$\Ens$ que satisfacen la condición
\ifthenelse{\boolean{orig}}
{subrayada más arriba.}
{destacada más arriba.}
\end{proposition}

En efecto, sea~$F_0$ un objeto del grupoide fundamental y sea
$\pi_0=\Aut(F_0)$. El funtor $\xi\mto\xi(F_0)$ es entonces una equivalencia
de la segunda categoría considerada en~\Ref{V.5.8} con la categoría
$\cal{C}(\pi_0)$, como se comprueba inmediatamente. Por otra parte, la
composición de este último con el funtor $X\mto E_X$ es la equivalencia
natural $\cal{C}\to\cal{C}(\pi_0)$. Se sigue que el propio funtor
$X\mto E_X$ es una equivalencia.

\begin{corollaire}
\label{V.5.9}
La categoría $\Prodash\cal{C}$
es canónicamente equivalente a la categoría de los funtores covariantes
$\xi$ del grupoide fundamental~$\Gamma$ en la categoría de los espacios
topológicos que satisfacen la condición siguiente: para todo objeto~$F$
de~$\Gamma$, $\xi(F)$ es un espacio compacto totalmente desconectado sobre
el que actúa el grupo topológico~$\pi_F$.
\end{corollaire}

También aquí, para verificar esta condición sobre~$\xi$, basta verificarla
para \emph{un}~$F$. La demostración es la misma que la de~\Ref{V.5.8}.

\begin{remarque}
\label{V.5.10}
Sea $(F_s)_{s\in S}$ una familia
\ifthenelse{\boolean{orig}}{de objeto}{de objetos}
del grupoide fundamental~$\Gamma$. Para $s,s'\in S$, pongamos
$$
\Hom(s,s')=\Hom(F_s,F_{s'})
$$
de modo que~$S$ se convierte a su vez en un grupoide conexo y la aplicación
$s\mto F_s$ en un funtor plenamente fiel de~$S$ en~$\Gamma$, que
denotaremos por~$f$. Al considerar entonces el funtor $X\mto E_X\circ f$
de~$\cal{C}$ en la categoría de funtores $\Hom(S,\Ens)$, se obtiene una
variante de~\Ref{V.5.8} (y de~\Ref{V.5.9}) con~$\Gamma$ reemplazado por~$S$.
El enunciado así obtenido se reduce al teorema~\Ref{V.4.1} cuando~$S$ se
reduce a un punto, y no es sino el propio enunciado~\Ref{V.5.8} si~$S$ es
el conjunto de los objetos de~$\Gamma$.
\end{remarque}

Utilizaremos~\Ref{V.5.9} para definir un proobjeto canónico de~$\cal{C}$.
Para ello consideramos el funtor de~$\Gamma$ en la categoría de los espacios
topológicos
\marginpar{132}
(e incluso de los grupos topológicos):
$$
f\colon F\mto\Aut(F)=\pi_F.
$$
Este funtor satisface la condición considerada en~\Ref{V.5.8}: el espacio
con operadores~$f(F)$ bajo~$\pi_F$ no es sino~$\pi_F$, considerado como
espacio con operadores bajo sí mismo mediante los automorfismos interiores.
Por tanto, el funtor~$f$ corresponde a un proobjeto de~$\cal{C}$
determinado salvo isomorfismo único, que es incluso un progrupo de
$\cal{C}$ y se llama el \emph{progrupo fundamental de~$\cal{C}$},
\index{progrupo fundamental de una categoría galoisiana|hyperpage}%
el cual desempeña el papel de un sistema local de grupos fundamentales. Es,
por tanto, un progrupo~$\Pi$ de~$\cal{C}$ definido por la condición de que
exista un isomorfismo funtorial en~$F$
$$
F(\Pi)\simeq\pi_F.
$$

Si~$X$ es un proobjeto cualquiera de~$\cal{C}$, se tiene un morfismo
canónico
$$
\Pi\times X\to X
$$
que hace de~$X$ un objeto con grupo de operadores izquierdos
\ifthenelse{\boolean{orig}}{$G$}{$\Pi$}
en~$\Prodash\cal{C}$. Para verlo basta observar que, para~$F$ variable, se
tiene una aplicación canónica
$$
\Pi(F)\times X(F)\to X(F)
$$
\ie
$$
\Aut(F)\times E_X(F)\to E_X(F),\quad\text{o}\quad
\pi_F\times F(X)\to F(X)
$$
que es funtorial en~$F$. También es funtorial en~$X$; por consiguiente, para
todo morfismo $X\to Y$ de proobjetos, el diagrama correspondiente
$$
\xymatrix{
\Pi\times X\ar[r]\ar[d]&X\ar[d]\\
\Pi\times Y\ar[r]&Y
}
$$
es conmutativo.

\begin{remarque}
\label{V.5.11}
No debe confundirse un proobjeto fundamental~$P$ (que no está provisto de
una estructura de grupo y es conexo) con el progrupo fundamental (que es un
pro\emph{grupo} y, en general, no es conexo). Más precisamente,
\ifthenelse{\boolean{orig}}{$G$}{$\Pi$}
es conexo si
\ifthenelse{\boolean{orig}}{\ignorespaces}{y}
solo si~$\pi_F$, actuando sobre sí mismo mediante automorfismos interiores,
actúa transitivamente, \ie si~$\pi$ se reduce al elemento neutro o,
equivalentemente, si~$\cal{C}$ es equivalente a la categoría de los
conjuntos finitos. Otra diferencia
\marginpar{133}
esencial es que
\ifthenelse{\boolean{orig}}{$G$}{$\Pi$}
está determinado salvo isomorfismo único, mientras que~$P$ solo está
determinado salvo isomorfismo no único.

Sea~$E$ un conjunto finito y consideremos el funtor constante sobre el
grupoide~$\Gamma$, de valor~$E$; en virtud de~\Ref{V.5.8}, este define un
objeto de~$\cal{C}$, denotado por~$E_{\cal{C}}$, que puede interpretarse
también como la suma de~$E$ ejemplares del objeto final~$e_{\cal{C}}$
de~$\cal{C}$. Puede considerarse~$E_{\cal{C}}$ como un funtor en~$E$, de la
categoría de los conjuntos finitos en la categoría~$\cal{C}$, y este funtor
es \emph{exacto}; por tanto, transforma grupos finitos en
$\cal{C}$-grupos,
\ifthenelse{\boolean{orig}}{etc...}{etc.}
Así pues, si~$X$ es un objeto de~$\cal{C}$ sobre el que el grupo finito~$G$
actúa por la derecha, puede considerarse~$X$ como un objeto de~$\cal{C}$
provisto del $\cal{C}$-grupo de operadores derechos~$G_{\cal{C}}$. Por
extensión de la terminología general relativa a objetos con grupos de
operadores que son $\cal{C}$-grupos, se dirá que~$X$ es \emph{formalmente
homogéneo principal}
\index{formalmente homogéneo principal (objeto)|hyperpage}%
bajo~$G$ si es formalmente homogéneo principal bajo~$G_{\cal{C}}$, \ie si
el morfismo canónico
$$
X\times G_{\cal{C}}\to X\times X
$$
deducido de la acción derecha de~$G_{\cal{C}}$ sobre~$X$, es un
isomorfismo. Se dice que~$X$ es \emph{homogéneo principal}
\index{homogéneo principal (objeto)|hyperpage}%
bajo~$G$ si lo es bajo~$G_{\cal{C}}$, \ie si lo es formalmente y, además,
el cociente $X/G=X/G_{\cal{C}}$ es~$e_{\cal{C}}$. Si se fija un funtor
fundamental, y por tanto una equivalencia de~$\cal{C}$ con una categoría
$\cal{C}(\pi)$, el objeto~$X$ corresponde a un conjunto sobre el que~$\pi$
actúa continuamente por la izquierda, sea~$E=F(X)$. Hacer actuar~$G$ por la
derecha sobre~$X$ equivale entonces a hacerlo actuar por la derecha sobre
el conjunto~$E$, de modo que las operaciones de~$G$ conmuten con las
de~$\pi$. Se comprueba de inmediato que~$X$ es homogéneo principal bajo~$G$
si y solo si el conjunto~$E$ es un espacio homogéneo principal
\ifthenelse{\boolean{orig}}
{bajo~$G$}
{bajo~$G$,}
\ie si y solo si~$G$ actúa sobre él de manera simplemente transitiva. (Por
otra parte, $X$ es formalmente homogéneo principal \sss\ 
$E$ es homogéneo principal \emph{o} vacío.) Comparando con~\Ref{V.5.3}, se
ve que, si~$X$ es homogéneo principal bajo~$G$ \emph{y conexo}, entonces el
homomorfismo dado de~$G$ en el grupo opuesto a~$\Aut(X)$ es un
\emph{isomorfismo}; además, para que un objeto~$X$
\ifthenelse{\boolean{orig}}{de~$C$}{de~$\cal{C}$}
sea conexo y homogéneo principal bajo el grupo opuesto a~$\Aut(X)$, es
necesario y suficiente, con las notaciones del \No~\Ref{V.4}, que sea
isomorfo a un~$P_i$ \emph{galoisiano}. En el caso modelo
$\cal{C}=\cal{C}(\pi)$, esto significa que~$X$ es isomorfo a un cociente
de~$\pi$ por un subgrupo normal.

Supongamos aún dado un funtor fundamental~$F$. Entonces, dar un
\marginpar{134}
objeto~$X$ homogéneo principal bajo un grupo finito~$G$ que actúa por la
derecha, junto con un punto $a\in F(X)$, equivale a dar un homomorfismo
de~$\pi$ en el grupo~$G$. En efecto, a tal homomorfismo se le hace
corresponder el conjunto~$E=G$, haciendo actuar~$\pi$ por la izquierda
mediante el homomorfismo dado $\pi\to G$ y las traslaciones izquierdas
de~$G$, y haciendo actuar~$G$ por la derecha mediante traslación derecha;
el punto marcado~$a$ de~$E$ es el elemento neutro de~$G$. Gracias a lo
anterior se obtiene así, de manera esencialmente única, todo triple
$(X,G,a)$ con las propiedades consideradas,
\ifthenelse{\boolean{orig}}
{puesto que un}
{puesto que un}
conjunto homogéneo principal bajo un grupo~$G$ y provisto de un punto
marcado se identifica con dicho grupo. De este modo se obtiene una
interpretación geométrica directa del funtor $G\mto\Hom(\pi,G)$ de la
categoría de los grupos finitos en~$\Ens$, funtor que es prorepresentable
mediante~$\pi$ y cuya consideración proporcionaría, por tanto, otra
construcción del grupo~$\pi$ asociado a~$F$.
\end{remarque}
